\documentclass[10pt,aspectratio=169]{beamer}
\usepackage[utf8]{inputenc}
\usepackage{listings}
\usepackage{hyperref}
\usepackage{xcolor}
\usepackage{ragged2e}
\usepackage{graphicx}
\usepackage{tikz}
\usetikzlibrary{positioning,calc,decorations.pathreplacing}

\usetheme{Dresden}
%\usetheme{Bergen}

\title{Programming Techniques 2026--2027}
\author{Hannu Parviainen}
\institute{Universidad de la Laguna}
\date{\today}

\lstset{
  language=fortran,
  basicstyle=\ttfamily\scriptsize,
  columns=fullflexible,
  keepspaces=true,
  backgroundcolor=\color{white},
  keywordstyle=\color{blue},
  commentstyle=\color{gray},
  stringstyle=\color{red},
  showstringspaces=false,
  breaklines=true,
  frame=lines
}

\subtitle{Lecture 5: Dynamic Memory and Pointers}

\begin{document}

\begin{frame}
  \titlepage
\end{frame}

\section{Dynamic Memory Allocation}

\begin{frame}
	\centering
	\Huge \textbf{Dynamic Memory Allocation}
\end{frame}


\begin{frame}[fragile]{Static memory allocation}
  \begin{columns}[T]
    \column{0.48\textwidth}
    \begin{block}{Concept}
      \begin{itemize}
        \item The array size is fixed at compile time.
        \item The memory is reserved when the program or procedure starts.
        \item Lifetime = the whole run of the program (or of the procedure).
        \item Efficient, but inflexible when the array size is not known in advance.
      \end{itemize}
    \end{block}
    \column{0.52\textwidth}
    \lstinputlisting[language=Fortran]{ex01.f90}
  \end{columns}
\end{frame}


\begin{frame}[fragile]{Dynamic memory allocation}
  \begin{columns}[T]
    \column{0.48\textwidth}
    \begin{block}{Concept}
      \begin{itemize}
        \item The size and shape are chosen at run time.
        \item The memory is allocated with \texttt{allocate}.
        \item It is released with \texttt{deallocate}, or automatically when the array goes out of scope.
        \item More flexible, at the cost of a little time for each allocation.
      \end{itemize}
    \end{block}
    \column{0.52\textwidth}
    \lstinputlisting[language=Fortran]{ex02.f90}
  \end{columns}
\end{frame}


\begin{frame}[fragile]{Allocatable Arrays in Fortran}
\begin{block}{Declaration}
\begin{lstlisting}[language=Fortran]
integer, dimension(:),   allocatable :: ages
real,    dimension(:,:), allocatable :: speed
\end{lstlisting}
 \end{block}
 \begin{block}{Allocation}
\begin{lstlisting}[language=Fortran]
integer :: isize, ierr
read *, isize

allocate(ages(isize))

allocate(speed(0:isize-1, 10), stat=ierr)
if (ierr /= 0) error stop 'speed: allocation failed'
\end{lstlisting}
 \end{block}
 \begin{block}{Note}
   Without \texttt{stat}, a failed allocation stops the program with an error message.
   With \texttt{stat}, the program continues, so you must handle the failure yourself.
 \end{block}
\end{frame}


\begin{frame}[fragile]{Automatic allocation on assignment}
  \begin{columns}[T]
    \begin{column}{0.45\textwidth}
      \begin{itemize}
        \item Since Fortran 2003, assigning to an allocatable array allocates it with the shape of the right-hand side.
        \item If the shapes differ, the array is reallocated. So \texttt{a = [a, x]} appends an element.
        \item Appending this way copies the whole array every time. That is fine for small arrays, but slow inside a long loop.
        \item \texttt{a(:) = ...} assigns to the existing elements and never reallocates.
      \end{itemize}
    \end{column}

    \begin{column}{0.55\textwidth}
      \begin{lstlisting}
program auto_alloc
  implicit none
  real, allocatable :: a(:)
  a = [1.0, 2.0, 3.0]     ! allocated with size 3
  a = [a, 4.0]            ! reallocated with size 4
  print *, size(a)        ! 4
  print *, a              ! 1.0  2.0  3.0  4.0
  a = a(2:3)              ! reallocated with size 2
  print *, size(a)        ! 2
  print *, a              ! 2.0  3.0
end program auto_alloc
      \end{lstlisting}
    \end{column}
  \end{columns}
\end{frame}


\begin{frame}[fragile]{Deallocation and Status}
 \begin{block}{Deallocate}
\begin{lstlisting}[language=Fortran]
if (allocated(ages)) deallocate(ages)
if (allocated(speed)) deallocate(speed)
\end{lstlisting}
 \end{block}
 \begin{block}{Guidelines}
   \begin{itemize}
     \item A local allocatable array is deallocated automatically when its procedure returns. Allocatable arrays cannot leak memory.
     \item Use \texttt{deallocate} to free memory early, for example a large work array.
     \item Allocating an array that is already allocated, or deallocating one that is not, is an error. Test with \texttt{allocated(array)}.
   \end{itemize}
 \end{block}
\end{frame}


\begin{frame}[fragile]{Allocatable strings}
  \begin{columns}[T]
    \begin{column}{0.45\textwidth}
      \begin{itemize}
        \item \texttt{character(len=:), allocatable} declares a string whose length is set when it is assigned.
        \item The string is reallocated whenever it is given a value of a different length.
        \item Compare with a fixed-length \texttt{character(len=10)}, which is padded with blanks and needs \texttt{trim}.
      \end{itemize}
    \end{column}

    \begin{column}{0.55\textwidth}
      \begin{lstlisting}
program strings
  implicit none
  character(len=:), allocatable :: name
  name = 'Vega'
  print *, name, len(name)
  name = name // ' (alpha Lyr)'
  print *, name, len(name)
end program strings
      \end{lstlisting}
      \begin{lstlisting}[language=bash]
$ ./a.out
 Vega           4
 Vega (alpha Lyr)          16
      \end{lstlisting}
    \end{column}
  \end{columns}
\end{frame}



\section{Pointers}

\begin{frame}
	\centering
	\Huge \textbf{Pointers}
\end{frame}

\begin{frame}{Fortran Pointers: Introduction}
  \begin{block}{What is a pointer?}
    \begin{itemize}
      \item A \textbf{pointer} is a variable that does not store a value directly,
            but instead \textbf{points to a memory location}.
      \item In Fortran, pointers can be associated with
            \begin{itemize}
              \item other variables,
              \item array sections,
              \item dynamically allocated memory.
            \end{itemize}
    \end{itemize}
  \end{block}

  \begin{block}{Why are they useful?}
    \begin{itemize}
      \item Allow multiple variables to \textbf{share the same data}.
      \item Needed for building \textbf{linked data structures}
            (lists, trees, graphs).
      \item Provide flexibility similar to references in modern languages.
    \end{itemize}
  \end{block}
\end{frame}


\begin{frame}{Python vs.~Fortran Variables}
  \begin{block}{Key difference}
    \begin{itemize}
      \item \textbf{Python variables} are \emph{references} to objects in memory.
        \begin{itemize}
          \item Assigning one variable to another does not copy the object.
          \item Example: \texttt{b = a} makes \texttt{b} point to the same object as \texttt{a}.
        \end{itemize}
      \item \textbf{Fortran variables} are normally \emph{values} stored directly.
        \begin{itemize}
          \item Assigning one variable to another copies the value.
          \item No implicit references like in Python.
        \end{itemize}
    \end{itemize}
  \end{block}

  \textbf{Python variables act like Fortran pointers by default}, while plain Fortran variables are independent copies.

\end{frame}


\begin{frame}[fragile]{Pointer Declaration}
  \begin{block}{Definition}
    \begin{itemize}
      \item A variable with the \texttt{pointer} attribute.
      \item Has static type, kind, and rank determined by its declaration.
    \end{itemize}
  \end{block}

  \begin{block}{Examples}
\begin{lstlisting}[language=Fortran]
real, pointer :: ptor => null()
real, dimension(:,:), pointer :: ptoa => null()
\end{lstlisting}
    \begin{itemize}
      \item \texttt{ptor} is a pointer to a scalar real target.
      \item \texttt{ptoa} is a pointer to a 2D array of reals.
      \item A new pointer points nowhere in particular: its status is \emph{undefined}. Initialise it with \texttt{=> null()}.
    \end{itemize}
  \end{block}
\end{frame}


\begin{frame}[fragile]{Target Declaration}

  \begin{block}{Targets must have the \texttt{target} attribute}
\begin{lstlisting}[language=Fortran]
real, target :: x, y
real, dimension(5,3), target :: a, b
real, dimension(3,5), target :: c
\end{lstlisting}
    \begin{itemize}
      \item \texttt{x} or \texttt{y} may become associated with \texttt{ptor}.
      \item \texttt{a}, \texttt{b}, or \texttt{c} may become associated with \texttt{ptoa}.
    \end{itemize}
  \end{block}
\end{frame}

\begin{frame}{Pointer Manipulation}
  \begin{block}{Operators}
    \begin{itemize}
      \item \texttt{=>} pointer assignment: alias pointer with a target.
      \item \texttt{=} normal assignment: assign value to space pointed at.
    \end{itemize}
    Pointer assignment makes the pointer and target reference the same space, while normal assignment changes the value.
  \end{block}
\end{frame}

\begin{frame}[fragile]{Pointer Assignment}
\begin{lstlisting}[language=Fortran]
real, target :: x, y
real, pointer :: ptor

x = 3.14159
ptor => y
ptor = x
\end{lstlisting}
  \begin{itemize}
    \item \texttt{x} and \texttt{ptor} have the same value.
    \item \texttt{ptor} is an alias for \texttt{y}, so the assignment sets \texttt{y = 3.14159}.
    \item Changing \texttt{x} later does not change \texttt{ptor} or \texttt{y}.
  \end{itemize}
\end{frame}

\begin{frame}[fragile]{Dynamic Targets}

  \begin{block}{Targets can be created dynamically by allocation}
\begin{lstlisting}[language=Fortran]
allocate(ptor, stat=ierr)
allocate(ptoa(n*n, 2*k-1), stat=ierr)
\end{lstlisting}
    \begin{itemize}
      \item First: allocates a scalar real as target of \texttt{ptor}.
      \item Second: allocates a rank-2 real array as target of \texttt{ptoa}.
      \item Allocating a pointer that is already associated is not an error. But if nothing else points to its old target, that memory is lost: a \textbf{memory leak}.
    \end{itemize}
  \end{block}
\end{frame}

\begin{frame}[fragile]{Association Status}
  \begin{block}{Test with \texttt{associated()}}
\begin{lstlisting}[language=Fortran]
associated(ptoa)        ! is it associated with any target?
associated(ptoa, arr)   ! is it associated with arr?
\end{lstlisting}
    \begin{itemize}
      \item Returns \texttt{.true.} if the pointer is associated.
      \item The second form tests association with a specific target.
      \item The test is only valid when the status of the pointer is defined. This is another reason to initialise pointers with \texttt{null()}.
    \end{itemize}
  \end{block}
\end{frame}

\begin{frame}[fragile]{Pointer Disassociation}
  \begin{block}{Nullification}
\begin{lstlisting}[language=Fortran]
nullify(ptor)       ! or: ptor => null()
\end{lstlisting}
    \begin{itemize}
      \item Breaks the pointer–target connection. The target itself is not changed.
    \end{itemize}
  \end{block}

  \begin{block}{Deallocation}
\begin{lstlisting}[language=Fortran]
deallocate(ptoa, stat=ierr)
\end{lstlisting}
    \begin{itemize}
      \item Breaks the connection and frees the target.
      \item Only for targets that were created with \texttt{allocate}. Deallocating a pointer to an ordinary variable is an error.
      \item Any other pointer to the same target is left \textbf{dangling}: it points to memory that no longer exists.
    \end{itemize}
  \end{block}
\end{frame}


\begin{frame}{Allocatable or pointer?}
  \begin{block}{Use allocatable arrays by default}
    They are freed automatically, so they cannot leak or dangle. The compiler also knows that they do not overlap other variables, which lets it produce faster code.
  \end{block}
  \begin{block}{Use pointers when you need what only they can do}
    Giving a second name to existing data, such as a section of a large array. Swapping large arrays without copying them. Building linked structures such as lists and trees.
  \end{block}
  \begin{block}{Pointer pitfalls}
    \begin{itemize}
      \item \textbf{Memory leak}: the last pointer to allocated memory is redirected or nullified before the memory is deallocated.
      \item \textbf{Dangling pointer}: the target is deallocated, or goes out of scope, while a pointer still points to it.
    \end{itemize}
  \end{block}
\end{frame}


\begin{frame}[fragile]{Practical Example: Swapping Arrays without Copying}
  \begin{columns}[T]
    \begin{column}{0.38\textwidth}
      \begin{itemize}
        \item Heat diffusion on a grid: each new value is the mean of its four neighbours.
        \item Every iteration needs the previous grid and the next one.
        \item Swapping the two pointers swaps their roles, and no data are copied.
        \item The loop ends when the largest change, \texttt{maxval(abs(...))}, is small enough. Here that takes 1898 iterations.
      \end{itemize}
    \end{column}

    \begin{column}{0.62\textwidth}
      \begin{lstlisting}[aboveskip=0pt, belowskip=0pt]
program jacobi
  implicit none
  integer, parameter :: n = 100
  real, target :: grid1(n,n), grid2(n,n)
  real, pointer :: old(:,:), new(:,:), tmp(:,:)
  integer :: iter
  grid1 = 0.0
  grid1(1,:) = 1.0        ! hot top edge
  grid2 = grid1
  old => grid1
  new => grid2
  do iter = 1, 100000
    new(2:n-1,2:n-1) = (old(1:n-2,2:n-1) + old(3:n,2:n-1) &
                      + old(2:n-1,1:n-2) + old(2:n-1,3:n)) / 4
    if (maxval(abs(new - old)) < 1.0e-4) exit
    tmp => old            ! swap the roles, copy nothing
    old => new
    new => tmp
  end do
  print *, 'Converged after', iter, 'iterations'
end program jacobi
      \end{lstlisting}
    \end{column}
  \end{columns}
\end{frame}


\section{Exercises}

\begin{frame}{Exercises: a growing array}
  \begin{block}{Tasks}
    \begin{enumerate}
      \item Write a program that reads real numbers from standard input until the input ends, and stores them in an allocatable array. Start from an empty array, \texttt{allocate(a(0))}, and append each new value with \texttt{a = [a, x]}.
      \item Print the number of values, their mean, and the values in reverse order.
      \item Test the program by redirecting a file to its standard input.
      \item What does the program do if the input is empty? Make it stop with a clear message.
      \item Appending copies the whole array each time. How could you read a very long file more efficiently?
    \end{enumerate}
  \end{block}
\end{frame}


\begin{frame}{Exercises: pointers}
  \begin{block}{Tasks}
    \begin{enumerate}
      \item Declare two real target arrays, \texttt{a} and \texttt{b}, with different values, and two pointers, \texttt{p} and \texttt{q}, initialised with \texttt{null()}.
      \item Point \texttt{p} at \texttt{a} and \texttt{q} at \texttt{b}. Change one element through \texttt{p} and print \texttt{a}. What happened?
      \item Swap the two pointers, so that \texttt{p} points at \texttt{b} and \texttt{q} at \texttt{a}, without copying any array.
      \item Confirm the result with \texttt{associated(p, b)} and \texttt{associated(q, a)}.
      \item Allocate a new target for \texttt{p} with \texttt{allocate}. Which of the pitfalls from the lecture can now happen, and how do you avoid them?
    \end{enumerate}
  \end{block}
\end{frame}

\begin{frame}{Exercises: sigma clipping}
  \begin{block}{Tasks}
    \begin{enumerate}
      \item Write a module with a function \texttt{clip(x, nsigma)} that returns the elements of the real array \texttt{x} that lie within \texttt{nsigma} standard deviations of its mean.
      \item The size of the result is only known at run time, so make the function result allocatable. Hint: count the elements to keep, allocate the result, and then copy them.
      \item Test it with \mbox{\texttt{[10.1, 9.8, 10.3, 9.9, 55.0, 10.0, 10.2]}} and \texttt{nsigma = 2.0}. Which value is removed?
      \item Store the result in an allocatable array in the main program. Do you need to allocate that array yourself?
      \item Repeat the clipping until no more values are removed, and print the mean of what is left.
    \end{enumerate}
  \end{block}
\end{frame}


\begin{frame}{Exercises: a histogram}
  \begin{block}{Tasks}
    \begin{enumerate}
      \item Write a program that reads integers from standard input until the input ends, and stores them in an allocatable array.
      \item Allocate a second array, \texttt{counts}, whose lower and upper bounds are the smallest and the largest value that was read.
      \item Count how many times each value occurs, so that \texttt{counts(v)} is the number of times the value \texttt{v} was read. Print every value with its count.
      \item Print the bounds and the size of \texttt{counts} with \texttt{lbound}, \texttt{ubound}, and \texttt{size}.
      \item What happens with empty input? Make the program stop with a clear message.
    \end{enumerate}
  \end{block}
\end{frame}


\begin{frame}{Exercises: a window into an image}
  \begin{block}{Tasks}
    \begin{enumerate}
      \item Read two integers, \texttt{nx} and \texttt{ny}, allocate a real array \texttt{image(nx,ny)} with the \texttt{target} attribute, and set it to zero.
      \item Point a two-dimensional pointer \texttt{window} at the central part of the image, \mbox{\texttt{image(nx/4+1:3*nx/4, ny/4+1:3*ny/4)}}. Print its shape and its lower bounds.
      \item Set \texttt{window = 1.0} and print the sum of the whole image. What does the result tell you?
      \item Point a one-dimensional pointer at every second element of the first column of the image, set those elements to 5.0, and print the first column.
      \item Deallocate the image. What does \texttt{associated(window)} return now, and can you trust it? What should you do with the pointers?
    \end{enumerate}
  \end{block}
\end{frame}



\end{document}
